\documentclass[12pt,a4paper]{article}
\usepackage[utf8]{inputenc}
\usepackage[T5]{fontenc}
\usepackage[vietnamese]{babel}
\usepackage[margin=1.3cm,top=1.0cm,bottom=1.05cm,includehead,headheight=18pt,headsep=8pt]{geometry}
\usepackage{amsmath,amssymb,array,tabularx,enumitem,fancyhdr,tikz,xcolor,needspace}
\usetikzlibrary{calc,angles,quotes,arrows.meta,babel,patterns}
\definecolor{ink}{RGB}{30,55,70}

\newcommand{\headerQuote}{%
  \ifcase\value{page}\relax
    \or Math gives us hope that every problem has a solution% Trang 1
    \or Math + You = Champion% Trang 2
    \or Where Frustration Becomes Inspiration% Trang 3
    \or Turning Confusion into Clarity% Trang 4
    \or Don't be an X, be a Y of mathematics% Trang 5
    \or Don't get mad. Learn to add% Trang 6
    \or The only way to learn mathematics is to do mathematics% Trang 7
    \or Small steps every day lead to big results% Trang 8
    \or Every expert was once a beginner% Trang 9
    \or Success is the sum of small efforts repeated daily% Trang 10
    \or Mathematics is the music of reason% Trang 11
    \or Believe you can and you're halfway there% Trang 12
    \else Math gives us hope that every problem has a solution%
  \fi
}
\pagestyle{fancy}\fancyhf{}
\lhead{\small\textit{\headerQuote}}
\rhead{\small\textit{Người biên soạn: Dương Gia Bảo}}
\rfoot{\small Trang \thepage}
\renewcommand{\headrulewidth}{0.4pt}
\setlength{\headheight}{18pt}
\setlength{\parindent}{0pt}
\setlength{\parskip}{0pt}
\setlength{\tabcolsep}{1.5pt}
\renewcommand{\arraystretch}{1.0}
\newcommand{\parttitle}[1]{\par\vspace{1.5pt}\needspace{2\baselineskip}{\large\bfseries\color{ink}#1}\par\nopagebreak\vspace{1pt}}
\newcommand{\four}[4]{\par\noindent\begin{tabularx}{\linewidth}{@{}*{4}{>{\raggedright\arraybackslash}X}@{}}\textbf{A.} #1&\textbf{B.} #2&\textbf{C.} #3&\textbf{D.} #4\end{tabularx}}
\newcommand{\two}[4]{\par\noindent\begin{tabularx}{\linewidth}{@{}*{2}{>{\raggedright\arraybackslash}X}@{}}\textbf{A.} #1&\textbf{B.} #2\\[4.5pt]\textbf{C.} #3&\textbf{D.} #4\end{tabularx}}
\newcommand{\one}[4]{\par\noindent\begin{tabularx}{\linewidth}{@{}>{\raggedright\arraybackslash}X@{}}\textbf{A.} #1\\[4pt]\textbf{B.} #2\\[4pt]\textbf{C.} #3\\[4pt]\textbf{D.} #4\end{tabularx}}
\newcommand{\answerbox}{\hfill\textit{Đáp số:}\ \fbox{\rule{0pt}{9pt}\hspace{2.2cm}}}
\setlist[enumerate,1]{label=\textbf{Câu \arabic*.},leftmargin=1.1cm,labelsep=.1cm,itemsep=0.1pt,topsep=0.1pt,parsep=0pt}
\setlist[enumerate,2]{label=\alph*),leftmargin=.55cm,itemsep=1.5pt,topsep=0.2pt,parsep=0pt}

\begin{document}
\noindent
\begin{minipage}[c]{0.38\textwidth}
    \centering
    \textbf{LÀM TOÁN ĐỂ HẠNH PHÚC}\\[-2pt]
    \rule{3.2cm}{0.6pt}\\[2pt]
    \textbf{ĐỀ CHÍNH THỨC}
\end{minipage}\hfill
\begin{minipage}[c]{0.35\textwidth}
    \centering
    \textbf{ĐỀ ÔN THI GIỮA KÌ I}\\
    \textbf{NĂM HỌC 2026 -- 2027}\\[2pt]
    \textbf{ĐỀ SỐ 34}\\
    \textbf{Môn: TOÁN -- Khối lớp 10}\\
    \textit{\small Thời gian làm bài: 90 phút (không kể thời gian phát đề)}
\end{minipage}

\vspace{0.01cm}
\noindent
\begin{minipage}[b]{0.78\textwidth}
    \textbf{Họ và tên:}\dotfill\textbf{SBD:}\makebox[3.5cm]{\dotfill}
\end{minipage}\hfill
\begin{minipage}[b]{0.20\textwidth}
    \hfill
    \setlength{\fboxrule}{0.8pt}%
    \fbox{\begin{minipage}{2.3cm}
        \centering
        \textbf{\small Mã đề thi}\\
        \textbf{0034}
    \end{minipage}}
\end{minipage}
\vspace{0.01cm}

\parttitle{PHẦN I. Trắc nghiệm nhiều lựa chọn}
\textit{Thí sinh trả lời từ câu 1 đến câu 12. Mỗi câu hỏi thí sinh chỉ chọn một phương án.}
\begin{enumerate}
\item Cho hai mệnh đề $P$ và $Q$. Phát biểu nào sau đây là đúng về mệnh đề kéo theo $P \Rightarrow Q$?
\two{$P$ là điều kiện đủ để có $Q$.}{$P$ là điều kiện cần để có $Q$.}{$Q$ là điều kiện đủ để có $P$.}{$P \Rightarrow Q$ sai khi $P$ sai và $Q$ đúng.}

\item Trong các mệnh đề sau, mệnh đề nào là mệnh đề \textbf{đúng}?
\four{$\exists x \in \mathbb{R},\ x^2 < 0$.}{$\forall n \in \mathbb{N},\ n^2 + n + 1$ lẻ.}{$\exists x \in \mathbb{Q},\ x^2 = 2$.}{$\forall x \in \mathbb{R},\ x^2 > x$.}

\item Cho hai tập hợp $A = (-3; 4]$ và $B = [1; 6)$. Tập hợp hiệu $A \setminus B$ là
\four{$(-3; 1]$.}{$[4; 6)$.}{$(-3; 1)$.}{$[1; 4]$.}

\item Cho hai tập hợp $A = [m - 1; m + 3]$ và $B = (-\infty; -2) \cup [5; +\infty)$ với $m$ là tham số thực. Có bao nhiêu số nguyên $m \in [-5; 10]$ để $A \cap B = \varnothing$?
\four{$1$.}{$2$.}{$4$.}{$3$.}

\item Hệ bất phương trình nào sau đây là hệ bất phương trình bậc nhất hai ẩn?
\two{$\begin{cases} 2x - 3y > -5 \\ x + 4y \le 1 \end{cases}$}{$\begin{cases} 2x^2 + y \le 3 \\ x - y > 0 \end{cases}$}{$\begin{cases} x - 3y + 2z \le 0 \\ 2x + y > 1 \end{cases}$}{$\begin{cases} xy + 1 > 0 \\ 3x - 2y \le 4 \end{cases}$}

\item Cho hệ bất phương trình $\begin{cases} 2x + y \le 4 \\ x - 2y \ge -2 \end{cases}$ ($y \ge 0$). Điểm nào sau đây thuộc miền nghiệm của hệ đã cho?
\four{$P(2; 3)$.}{$M(1; 1)$.}{$Q(3; 0)$.}{$N(0; 2)$.}

\item Nửa mặt phẳng (kể cả bờ $d: 3x - 2y = 6$) chứa điểm $M(1; -1)$ là miền nghiệm của bất phương trình nào sau đây?
\four{$3x - 2y \ge 6$.}{$3x - 2y > 6$.}{$3x - 2y \le 6$.}{$3x - 2y < 6$.}

\item Giá trị lớn nhất của biểu thức $F(x, y) = 2x + 3y$ trên miền $\begin{cases} x + y \le 8 \\ y \le 3x \end{cases}$ ($0 \le y \le 6, x \ge 0$) bằng
\four{$16$.}{$18$.}{$20$.}{$22$.}

\item Cho tam giác $ABC$ có bán kính đường tròn ngoại tiếp $R$. Khẳng định nào sau đây là đúng?
\four{$\dfrac{a}{\sin A} = 2R$.}{$\dfrac{a}{\sin A} = R$.}{$\dfrac{a}{\sin A} = 4R$.}{$\dfrac{a}{\cos A} = 2R$.}

\item Giá trị của biểu thức $M = \sin 150^\circ + \cos 120^\circ$ bằng
\four{$1$.}{$0$.}{$-1$.}{$\dfrac{\sqrt{3}}{2}$.}

\item Cho tam giác $ABC$ có độ dài hai cạnh $AB = 4, AC = 6$ và góc $\widehat{A} = 120^\circ$. Độ dài cạnh $BC$ bằng
\four{$2\sqrt{7}$.}{$\sqrt{28}$.}{$2\sqrt{19}$.}{$2\sqrt{13}$.}

\item Cho góc nhọn $\alpha$ thoả mãn $\sin \alpha = \dfrac{3}{5}$. Giá trị của biểu thức $P = \dfrac{2\cos \alpha - \tan \alpha}{1 + \sin \alpha}$ bằng
\four{$\dfrac{7}{32}$.}{$\dfrac{15}{32}$.}{$\dfrac{19}{32}$.}{$\dfrac{17}{32}$.}
\end{enumerate}
\newpage

\parttitle{PHẦN II. Trắc nghiệm đúng/sai}
\textit{Thí sinh trả lời từ câu 1 đến câu 2. Trong mỗi ý a), b), c), d) ở mỗi câu, thí sinh chọn đúng hoặc sai.}
\begin{enumerate}
\item Cho hai tập hợp $A = \{x \in \mathbb{N} \mid (2x^2 - 5x + 2)(x^2 - 9) = 0\}$, $B = \{x \in \mathbb{Z} \mid -2 \le x \le 3\}$ và tập hợp $C = [m - 2; m + 1]$ (với $m$ là tham số thực).
\begin{enumerate}
    \item Số phần tử của tập hợp $A$ bằng $2$.
    \item Tập hợp $B$ có tất cả $64$ tập hợp con.
    \item Tập hợp hiệu $A \setminus B = \{2; 3\}$.
    \item Có đúng $5$ giá trị nguyên của tham số $m$ để $A \cap C \ne \varnothing$.
\end{enumerate}

\item Một doanh nghiệp kinh doanh hai loại mặt hàng $X$ và $Y$. Giả sử số tấn mặt hàng $X$ và $Y$ bán ra trong một tuần lần lượt là $x$ và $y$ ($x \ge 0, y \ge 0$). Cặp số $(x; y)$ phải thoả mãn hệ bất phương trình:
\[
\begin{cases}
x + y \le 10 \\
2x + y \le 14 \\
x \ge 0,\ y \ge 0
\end{cases}
\]
Lợi nhuận thu được khi bán $x$ tấn hàng $X$ và $y$ tấn hàng $Y$ được cho bởi hàm số $F(x, y) = 3x + 2y$ (triệu đồng).
\begin{enumerate}
    \item Cặp số $(x; y) = (2; 3)$ là một nghiệm của hệ bất phương trình đã cho.
    \item Miền nghiệm của hệ bất phương trình là miền tứ giác có toạ độ các đỉnh là $(0; 0)$, $(7; 0)$, $(4; 6)$ và $(0; 10)$.
    \item Diện tích miền nghiệm của hệ bất phương trình bằng $45$.
    \item Doanh nghiệp có thể thu được lợi nhuận lớn nhất trong một tuần là $24$ triệu đồng.
\end{enumerate}
\end{enumerate}

\parttitle{PHẦN III. Trắc nghiệm trả lời ngắn}
\textit{Thí sinh trả lời từ câu 1 đến câu 4. Điền kết quả vào ô tương ứng.}
\begin{enumerate}
\item Cho miền tứ giác $OABC$ (kể cả biên) với toạ độ các đỉnh là gốc toạ độ $O(0; 0)$, $A(5; 0)$, $B(3; 4)$ và $C(0; 3)$ như hình vẽ.

\begin{center}
\vspace{-3pt}
\begin{tikzpicture}[scale=0.72, >=Stealth]
    \draw[line width=1.1pt, ->] (-0.8,0) -- (6.2,0) node[below=2pt, font=\small\bfseries] {$x$};
    \draw[line width=1.1pt, ->] (0,-0.8) -- (0,5.0) node[left=2pt, font=\small\bfseries] {$y$};
    \node[below left=2pt, font=\small\bfseries] at (0,0) {$O$};
    
    \fill[pattern=north east lines, pattern color=blue!50] (0,0) -- (5,0) -- (3,4) -- (0,3) -- cycle;
    \draw[line width=1.3pt, ink] (0,0) -- (5,0) -- (3,4) -- (0,3) -- cycle;
    
    \filldraw[black] (5,0) circle (2.4pt) node[below=2pt, font=\small\bfseries] {$A(5;0)$};
    \filldraw[black] (3,4) circle (2.4pt) node[above right=1pt, font=\small\bfseries] {$B(3;4)$};
    \filldraw[black] (0,3) circle (2.4pt) node[left=2pt, font=\small\bfseries] {$C(0;3)$};
    
    \draw[dashed, thin] (3,0) -- (3,4) -- (0,4);
    \node[below=2pt, font=\small] at (3,0) {$3$};
    \node[left=2pt, font=\small] at (0,4) {$4$};
\end{tikzpicture}
\vspace{-5pt}
\end{center}

Tìm giá trị lớn nhất của biểu thức $F(x, y) = 3x + 2y$ với $(x; y)$ thuộc miền tứ giác $OABC$.
\answerbox

\item Để đo chiều cao $CD$ của một ngọn tháp ven sông mà không thể đến chân tháp, người ta cắm hai cọc thẳng đứng tại hai điểm $A$ và $B$ trên mặt đất cách nhau một khoảng $AB = 40\text{ m}$ (sao cho $A, B, C$ thẳng hàng theo thứ tự đó). Bằng giác kế, người ta đo được góc nâng $\widehat{CAD} = 30^\circ$ và $\widehat{CBD} = 45^\circ$. Giả sử chiều cao của giác kế không đáng kể. Chiều cao $CD$ của ngọn tháp bằng bao nhiêu mét? (Làm tròn kết quả đến hàng đơn vị).
\answerbox
\end{enumerate}
\newpage

\begin{enumerate}[resume]
\item Trong đợt khảo sát sở thích thể thao của một khối lớp 10 gồm $120$ học sinh, có $65$ học sinh thích bóng đá, $50$ học sinh thích bóng rổ, $42$ học sinh thích bơi lội; có $25$ học sinh thích cả bóng đá và bóng rổ, $18$ học sinh thích cả bóng đá và bơi lội, $15$ học sinh thích cả bóng rổ và bơi lội. Biết rằng có $8$ học sinh thích cả ba môn thể thao trên. Hỏi có bao nhiêu học sinh trong khối 10 không thích môn nào trong ba môn thể thao trên?
\answerbox

\item Cho hai tập hợp khác rỗng $A = [m - 2; 2m + 5)$ và $B = (-\infty; -3) \cup [7; +\infty)$ với $m$ là tham số thực. Có bao nhiêu giá trị nguyên của $m$ để tập hợp $A \setminus B$ chứa đúng $4$ số nguyên?
\answerbox
\end{enumerate}

\parttitle{PHẦN IV. Tự luận}
\textit{Thí sinh trình bày chi tiết lời giải từ câu 1 đến câu 4.}
\begin{enumerate}
\item Để đo chiều cao $CD$ của một ngọn núi đá mà không thể đến chân núi, hai quan sát viên tại hai vị trí $A$ và $B$ trên mặt đất cách nhau $AB = 250\text{ m}$ cùng hướng ống ngắm về đỉnh núi $D$. Góc nâng nhìn từ $A$ lên $D$ là $32^\circ$, góc nâng nhìn từ $B$ lên $D$ là $40^\circ$. Biết rằng góc nhìn trên mặt đất $\widehat{ACB} = 75^\circ$ với $C$ là hình chiếu vuông góc của đỉnh $D$ lên mặt phẳng nằm ngang chứa $A$ và $B$.
\begin{enumerate}
    \item Biểu diễn khoảng cách $AC$ và $BC$ theo chiều cao $h = CD$ của ngọn núi.
    \item Tính chiều cao $CD$ của ngọn núi (làm tròn kết quả đến hàng phần mười theo đơn vị mét).
\end{enumerate}

\item Cho hai tập hợp khác rỗng $A = (-3; 4]$ và $B = [m - 1; 2m + 3)$ với $m$ là tham số thực.
\begin{enumerate}
    \item Khi $m = 1$, hãy xác định các tập hợp $A \cap B$, $A \cup B$ và $A \setminus B$.
    \item Tìm tất cả các giá trị thực của tham số $m$ để tập hợp $B$ là tập con của tập hợp $A$ ($B \subseteq A$).
\end{enumerate}

\item Một cơ sở sản xuất thủ công mỹ nghệ dự định sản xuất hai loại sản phẩm gốm: Bình hoa gốm (loại I) và Chậu cảnh gốm (loại II).
\begin{itemize}[leftmargin=1.5em, itemsep=0.1pt]
    \item Để sản xuất $1$ sản phẩm loại I cần $2\text{ kg}$ đất sét và $3$ giờ công lao động.
    \item Để sản xuất $1$ sản phẩm loại II cần $4\text{ kg}$ đất sét và $2$ giờ công lao động.
\end{itemize}
Cơ sở hiện có tối đa $40\text{ kg}$ đất sét và quỹ thời gian tối đa là $36$ giờ làm việc. Sau khi bán, mỗi sản phẩm loại I mang lại tiền lãi $150$ nghìn đồng, mỗi sản phẩm loại II mang lại tiền lãi $180$ nghìn đồng.
\begin{enumerate}
    \item Gọi $x, y$ lần lượt là số sản phẩm loại I và loại II mà cơ sở sản xuất trong một đợt ($x \ge 0, y \ge 0$). Hãy lập hệ bất phương trình bậc nhất hai ẩn thể hiện các điều kiện ràng buộc trên và biểu diễn miền nghiệm của hệ trên mặt phẳng toạ độ $Oxy$.
    \item Cơ sở nên sản xuất bao nhiêu sản phẩm mỗi loại để thu được tổng tiền lãi lớn nhất? Tổng số tiền lãi lớn nhất đó bằng bao nhiêu?
\end{enumerate}

\item Một bờ sông thẳng có hai ngôi làng $A$ và $B$ cùng nằm về một phía của bờ sông. Khoảng cách từ làng $A$ đến bờ sông là $AH = 3\text{ km}$, khoảng cách từ làng $B$ đến bờ sông là $BK = 5\text{ km}$. Khoảng cách giữa hai hình chiếu $H$ và $K$ trên bờ sông là $HK = 6\text{ km}$. Người ta muốn xây dựng một trạm bơm nước $M$ trên đoạn bờ sông $HK$ để dẫn nước về hai làng.
\begin{enumerate}
    \item Giả sử trạm bơm $M$ được đặt tại trung điểm của đoạn thẳng $HK$. Tính tổng độ dài đường ống dẫn nước $MA + MB$ (làm tròn kết quả đến hàng phần trăm theo đơn vị kilômét).
    \item Xác định vị trí điểm $M$ trên đoạn thẳng $HK$ để tổng chiều dài đường ống dẫn nước $MA + MB$ là ngắn nhất. Chiều dài ngắn nhất đó bằng bao nhiêu kilômét?
\end{enumerate}
\end{enumerate}

\newpage

% ==================== ĐÁP ÁN VÀ LỜI GIẢI CHI TIẾT ====================
\parttitle{BẢNG ĐÁP ÁN TỔNG HỢP}

\vspace{0.2cm}
\noindent\textbf{PHẦN I. Trắc nghiệm nhiều lựa chọn (12 câu)}
\vspace{0.15cm}

\begin{center}
\begin{tabular}{|c|c|c|c|c|c|c|c|c|c|c|c|}
\hline
\textbf{1} & \textbf{2} & \textbf{3} & \textbf{4} & \textbf{5} & \textbf{6} & \textbf{7} & \textbf{8} & \textbf{9} & \textbf{10} & \textbf{11} & \textbf{12} \\
\hline
\textbf{A} & \textbf{B} & \textbf{C} & \textbf{D} & \textbf{A} & \textbf{B} & \textbf{C} & \textbf{D} & \textbf{A} & \textbf{B} & \textbf{C} & \textbf{D} \\
\hline
\end{tabular}
\end{center}

\vspace{0.3cm}
\noindent\textbf{PHẦN II. Trắc nghiệm đúng/sai (2 câu)}
\vspace{0.15cm}

\begin{center}
\begin{tabular}{|c|c|c|c|c|}
\hline
\textbf{Câu} & \textbf{Lệnh hỏi a)} & \textbf{Lệnh hỏi b)} & \textbf{Lệnh hỏi c)} & \textbf{Lệnh hỏi d)} \\
\hline
\textbf{Câu 1} & \textbf{Đúng} & \textbf{Đúng} & \textbf{Sai} & \textbf{Đúng} \\
\hline
\textbf{Câu 2} & \textbf{Đúng} & \textbf{Đúng} & \textbf{Sai} & \textbf{Đúng} \\
\hline
\end{tabular}
\end{center}

\vspace{0.3cm}
\noindent\textbf{PHẦN III. Trắc nghiệm trả lời ngắn (4 câu)}
\vspace{0.15cm}

\begin{center}
\begin{tabular}{|c|c|c|c|c|}
\hline
\textbf{Câu} & \textbf{Câu 1} & \textbf{Câu 2} & \textbf{Câu 3} & \textbf{Câu 4} \\
\hline
\textbf{Đáp số} & $17$ & $55$ & $13$ & $2$ \\
\hline
\end{tabular}
\end{center}

\vspace{0.4cm}
\parttitle{PHẦN IV. ĐÁP SỐ TỰ LUẬN}
\vspace{0.15cm}

\begin{enumerate}[leftmargin=1.5em, itemsep=6pt]
    \item[\textbf{Câu 1.}] 
    \begin{itemize}[leftmargin=1.2em, itemsep=2pt]
        \item[a)] $AC = h \cot 32^\circ \approx 1{,}6003\, h$, \quad $BC = h \cot 40^\circ \approx 1{,}1918\, h$.
        \item[b)] $CD \approx 144{,}5\text{ m}$.
    \end{itemize}

    \item[\textbf{Câu 2.}] 
    \begin{itemize}[leftmargin=1.2em, itemsep=2pt]
        \item[a)] $A \cap B = [0; 4]$, \quad $A \cup B = (-3; 5)$, \quad $A \setminus B = (-3; 0)$.
        \item[b)] $-2 < m \le \dfrac{1}{2}$.
    \end{itemize}

    \item[\textbf{Câu 3.}] 
    \begin{itemize}[leftmargin=1.2em, itemsep=2pt]
        \item[a)] Hệ bất phương trình mô hình hóa nguyên liệu đất sét và công lao động:
        \[
        \begin{cases}
        x \ge 0,\ y \ge 0 \\
        x + 2y \le 20 \\
        3x + 2y \le 36
        \end{cases}
        \]
        \item[b)] Cơ sở nên sản xuất $8$ sản phẩm loại I và $6$ sản phẩm loại II để thu được tiền lãi lớn nhất là $2280$ nghìn đồng ($2{,}28$ triệu đồng).
    \end{itemize}

    \item[\textbf{Câu 4.}] 
    \begin{itemize}[leftmargin=1.2em, itemsep=2pt]
        \item[a)] $MA + MB = 3\sqrt{2} + \sqrt{34} \approx 10{,}07\text{ km}$.
        \item[b)] $(MA + MB)_{\min} = 10\text{ km}$ khi trạm bơm $M$ cách $H$ một khoảng $HM = 2{,}25\text{ km}$.
    \end{itemize}
\end{enumerate}

\newpage

% ==================== HÌNH VẼ MIỀN NGHIỆM HỆ BẤT PHƯƠNG TRÌNH ====================
\parttitle{HÌNH VẼ MIỀN NGHIỆM HỆ BẤT PHƯƠNG TRÌNH}

\vspace{0.2cm}
\noindent\textbf{1. Miền nghiệm Câu 1 -- Phần III (Tối ưu biểu thức $F(x, y) = 3x + 2y$ trên miền tứ giác $OABC$):}

\begin{center}
\begin{tikzpicture}[x=1.35cm, y=1.35cm]
    % Lưới toạ độ
    \draw[help lines, color=gray!25, step=1] (-1,-1) grid (7,6);
    
    % Trục toạ độ
    \draw[-{Latex[length=3mm]}, line width=1.1pt] (-1.2,0) -- (7.2,0) node[below=2pt, font=\normalsize\bfseries] {$x$};
    \draw[-{Latex[length=3mm]}, line width=1.1pt] (0,-1.2) -- (0,5.8) node[left=2pt, font=\normalsize\bfseries] {$y$};
    \node[below left=2pt, font=\normalsize\bfseries] at (0,0) {$O$};
    
    % Tô màu miền nghiệm OABC
    \fill[ink!20] (0,0) -- (5,0) -- (3,4) -- (0,3) -- cycle;
    \draw[line width=1.4pt, ink] (0,0) -- (5,0) -- (3,4) -- (0,3) -- cycle;
    
    % Các đường thẳng biên
    % d_AB qua A(5,0) và B(3,4): pt 2x + y = 10
    \draw[line width=1.1pt, blue!80!black] (2.2,5.6) -- (5.6,-1.2) 
        node[pos=0.08, above right, font=\small\bfseries] {$d_1: 2x + y = 10$};
        
    % d_BC qua B(3,4) và C(0,3): pt x - 3y = -9 <=> x - 3y + 9 = 0
    \draw[line width=1.1pt, red!80!black] (-0.8,2.73) -- (6.5,5.17) 
        node[pos=0.92, above left, font=\small\bfseries] {$d_2: x - 3y = -9$};

    % Các đường gióng toạ độ
    \draw[dashed, line width=0.8pt, gray!80] (3,0) node[below=2pt, font=\small\bfseries] {$3$} -- (3,4) -- (0,4) node[left=2pt, font=\small\bfseries] {$4$};
    \node[below=2pt, font=\small\bfseries] at (5,0) {$5$};
    \node[left=2pt, font=\small\bfseries] at (0,3) {$3$};

    % Các điểm đỉnh
    \fill[black] (0,0) circle (2.5pt);
    \fill[black] (5,0) circle (2.5pt) node[above right=1pt, font=\normalsize\bfseries] {$A(5;0)$};
    \fill[black] (3,4) circle (2.5pt) node[above right=1pt, font=\normalsize\bfseries] {$B(3;4)$};
    \fill[black] (0,3) circle (2.5pt) node[above right=1pt, font=\normalsize\bfseries] {$C(0;3)$};

    % Nhãn miền nghiệm
    \node[font=\bfseries\color{ink}] at (2.0,1.8) {Miền tứ giác $OABC$};
\end{tikzpicture}
\end{center}

\vspace{0.2cm}
\begin{itemize}[leftmargin=1.5em, itemsep=3pt]
    \item Miền nghiệm là tứ giác $OABC$ (kể cả biên) với toạ độ các đỉnh:
    \[
    O(0;0), \quad A(5;0), \quad B(3;4), \quad C(0;3).
    \]
    \item Hàm mục tiêu: $F(x, y) = 3x + 2y$.
    \item Giá trị hàm mục tiêu tại các đỉnh của tứ giác:
    \[
    F(O) = 3(0) + 2(0) = 0; \quad F(A) = 3(5) + 2(0) = 15;
    \]
    \[
    F(B) = 3(3) + 2(4) = 17; \quad F(C) = 3(0) + 2(3) = 6.
    \]
    \item \textbf{Kết luận:} Giá trị lớn nhất của $F(x, y)$ trên miền tứ giác $OABC$ là $\mathbf{17}$ (đạt tại đỉnh $B(3;4)$).
\end{itemize}

\newpage

% ==================== MIỀN NGHIỆM CÂU 3 PHẦN IV ====================
\parttitle{HÌNH VẼ MIỀN NGHIỆM HỆ BẤT PHƯƠNG TRÌNH (Tiếp theo)}

\vspace{0.2cm}
\noindent\textbf{2. Miền nghiệm Câu 3 -- Phần IV (Bài toán quy hoạch sản xuất gốm sứ thủ công mỹ nghệ):}
\[
\text{Hệ bất phương trình ràng buộc: }
\begin{cases}
x \ge 0,\ y \ge 0 \\
x + 2y \le 20 \\
3x + 2y \le 36
\end{cases}
\]

\begin{center}
\begin{tikzpicture}[x=0.72cm, y=0.72cm]
    % Lưới toạ độ
    \draw[help lines, color=gray!25, step=2] (-1,-1) grid (16,13);
    
    % Trục tọa độ
    \draw[-{Latex[length=3mm]}, line width=1.1pt] (-1.2,0) -- (16.5,0) node[below=2pt, font=\normalsize\bfseries] {$x$};
    \draw[-{Latex[length=3mm]}, line width=1.1pt] (0,-1.2) -- (0,13.5) node[left=2pt, font=\normalsize\bfseries] {$y$};
    \node[below left=2pt, font=\normalsize\bfseries] at (0,0) {$O$};

    % Tô màu miền nghiệm OABC
    \fill[ink!20] (0,0) -- (12,0) -- (8,6) -- (0,10) -- cycle;
    \draw[line width=1.4pt, ink] (0,0) -- (12,0) -- (8,6) -- (0,10) -- cycle;

    % d1: 3x + 2y = 36
    \draw[line width=1.1pt, blue!80!black] (4,12) -- (13.33,-2) 
        node[pos=0.08, above right, font=\small\bfseries] {$d_1: 3x + 2y = 36$};

    % d2: x + 2y = 20
    \draw[line width=1.1pt, red!80!black] (-1,10.5) -- (15.5,2.25) 
        node[pos=0.92, above right, font=\small\bfseries] {$d_2: x + 2y = 20$};

    % Đường gióng toạ độ
    \draw[dashed, line width=0.8pt, gray!80] (8,0) node[below=2pt, font=\small\bfseries] {$8$} -- (8,6) -- (0,6) node[left=2pt, font=\small\bfseries] {$6$};
    \node[below=2pt, font=\small\bfseries] at (12,0) {$12$};
    \node[left=2pt, font=\small\bfseries] at (0,10) {$10$};

    % Các đỉnh
    \fill[black] (0,0) circle (2.5pt);
    \fill[black] (12,0) circle (2.5pt) node[above right=1pt, font=\normalsize\bfseries] {$A(12;0)$};
    \fill[black] (8,6) circle (2.5pt) node[above right=1pt, font=\normalsize\bfseries] {$B(8;6)$};
    \fill[black] (0,10) circle (2.5pt) node[above right=1pt, font=\normalsize\bfseries] {$C(0;10)$};

    % Nhãn miền nghiệm
    \node[font=\bfseries\color{ink}] at (4.5,4.0) {Miền nghiệm $OABC$};
\end{tikzpicture}
\end{center}

\vspace{0.2cm}
\begin{itemize}[leftmargin=1.5em, itemsep=3pt]
    \item Miền nghiệm của hệ là tứ giác $OABC$ (kể cả biên) với toạ độ các đỉnh:
    \[
    O(0;0), \quad A(12;0), \quad B(8;6), \quad C(0;10).
    \]
    \item Hàm tiền lãi: $L(x, y) = 150x + 180y$ (nghìn đồng).
    \item Giá trị hàm tiền lãi tại các đỉnh:
    \[
    L(O) = 150(0) + 180(0) = 0 \text{ (nghìn đồng)};
    \]
    \[
    L(A) = 150(12) + 180(0) = 1800 \text{ (nghìn đồng)} = 1{,}8 \text{ triệu đồng};
    \]
    \[
    L(B) = 150(8) + 180(6) = 1200 + 1080 = 2280 \text{ (nghìn đồng)} = 2{,}28 \text{ triệu đồng};
    \]
    \[
    L(C) = 150(0) + 180(10) = 1800 \text{ (nghìn đồng)} = 1{,}8 \text{ triệu đồng}.
    \]
    \item \textbf{Kết luận:} Tiền lãi lớn nhất là $\mathbf{2280}$ nghìn đồng ($=\mathbf{2{,}28}$ triệu đồng), đạt được khi sản xuất $\mathbf{8}$ sản phẩm loại I và $\mathbf{6}$ sản phẩm loại II.
\end{itemize}

\end{document}
